\section{Background and Related Work}
\label{sec:background}

Blockchain fraud detection has progressed through three intellectual
generations.  The first applied hand-crafted heuristics and tabular
classifiers to transaction statistics~\cite{ostapowicz2019,blanco2022,
ashfaq2022,elmougy2023}.  The second exploited the graph structure of
the blockchain, treating addresses as nodes and transactions as
edges~\cite{kim2023,zhou2021,ctgcn2024}.  The third, now emerging,
combines multiple modalities to capture both structural and semantic
dimensions of fraud.  We survey each generation, identify the residual
gaps, and position our contribution.

\subsection{Graph-Centric Detection: Behavioural Forensics at Scale}

The foundational insight of graph-centric detection is that illicit
activity manifests as detectable anomalies in the topology of the
transaction network.  Weber \emph{et al.}~\cite{weber2019} established
this on the Elliptic Bitcoin dataset with graph convolutional networks,
showing that local topological features---degree, clustering
coefficient, temporal velocity---distinguish laundering rings from
benign transfers.  Huang \emph{et al.}~\cite{huang2024peae} scaled this
to Ethereum with PEAE-GNN, constructing augmentation ego-graphs around
known accounts to make GNN-based detection tractable at chain
scale.

Subsequent work addressed the two most pressing practical
limitations.  Chen \emph{et al.}~\cite{chen2024dahgnn} tackled class
imbalance---fraud addresses are rare even in seeded datasets---with
DA-HGNN, which combines data augmentation with temporal and structural
learning.  Li \emph{et al.}~\cite{li2025global,li2022ttagn} added
global context by introducing Transformer-style attention over transaction
graphs, enabling the model to detect distributed money-laundering schemes
that span multiple hops.  Ouyang \emph{et al.}~\cite{ouyang2024bitchetg}
extended graph analysis to subgraph-level contrastive learning for
identifying \emph{coordinated} laundering groups, while Kanezashi
\emph{et al.}~\cite{kanezashi2022hetero} showed that heterogeneous
graphs---which encode different node and edge types explicitly---
consistently outperform homogeneous alternatives on Ethereum phishing.

\textbf{Residual gap.}
Graph-centric methods are blind to what contracts \emph{do}.
Two contracts that appear structurally indistinguishable---similar
degrees, similar value flows, similar counterparty sets---can differ
categorically in intent: one may contain a hidden backdoor that drains
every depositor's balance on a single owner call.
This flaw is undetectable without analysing the source code.

\subsection{Code-Centric Detection: Semantic Vulnerability Analysis}

A parallel tradition treats smart contracts as executable artefacts and
analyses their bytecode or source for vulnerabilities.
Luu \emph{et al.}~\cite{luu2016oyente} introduced Oyente, a symbolic
execution engine that identifies reentrancy and transaction-ordering
vulnerabilities.  Tsankov \emph{et al.}~\cite{tsankov2018securify}
formalised compliance checking with Securify.  Qian \emph{et
al.}~\cite{qian2020reentrancy} showed that reentrancy can be learned
directly from opcode sequences with a bidirectional LSTM, eliminating
hand-crafted rules.  Torres \emph{et al.}~\cite{torres2019honeypots}
specialised in honeypot detection.  Zhang \emph{et
al.}~\cite{zhang2020txspector} replayed historical transaction traces
for forensic post-mortem analysis.  Sun \emph{et
al.}~\cite{sun2024gptscan} hybridised large language models with static
analysis in GPTScan, reducing false positives for logic vulnerabilities.
Kong \emph{et al.}~\cite{kong2025uechecker} applied graph learning to
call graphs within contracts, detecting unchecked external-call patterns.

\textbf{Residual gap.}
Code-centric methods are blind to the network context in which a
contract operates.  A contract whose source code is a clean, verified
fork of a legitimate DeFi protocol can still be a fraud instrument if
its transaction neighbourhood connects it to known exploit infrastructure.
Source analysis cannot detect this.

\subsection{Hybrid Detection: Towards Holistic Fusion}

The recognition that neither structural nor semantic analysis alone is
sufficient has driven the emergence of hybrid approaches.
Hu \emph{et al.}~\cite{hu2023bert4eth} introduced BERT4ETH, pre-training
a Transformer encoder over Ethereum transaction sequences treated as
sentences.  This captures behavioural patterns but remains limited to
transaction metadata---it does not ingest contract source code.
Liang \emph{et al.}~\cite{liang2025ponziguard} proposed PonziGuard,
constructing Contract Runtime Behaviour Graphs that couple bytecode
call sequences with fund-flow dynamics, achieving strong Ponzi detection
but targeting a narrow fraud subtype.
Sun \emph{et al.}~\cite{sun2025tlmg4eth} developed TLMG4Eth, one of
the first systems to jointly learn from transaction language and graph
representations, reporting fusion gains over isolated modalities.
Wu \emph{et al.}~\cite{wu2024tokenscout} extended graph evolution with
token attributes for early scam-token detection.
Galletta and Pinelli~\cite{galletta2024ponzi} emphasised
interpretability, using XAI techniques to show that features from both
transaction and code domains make fraud classifiers more transparent.

\textbf{Residual gap.}
Existing hybrid methods share three limitations that our work addresses.
First, none seeds their dataset from \emph{victim contracts with verified
source code}; most rely on attacker-wallet labels from Forta or
Ethereum phishing benchmarks, which provide no source for the LLM
encoder to analyse.
Second, no published system reports a module-level ablation that
compares alternative graph encoders (GCN, GAT) and fusion mechanisms
(concatenation vs.\ attention) on the same split.
Third, reproducibility varies: several systems were developed in private
notebooks on private datasets, making independent verification impossible.

\subsection{Positioning This Work}

Our system sits within the hybrid stream.  Relative to TLMG4Eth and
PonziGuard, it is intentionally minimal---two GraphSAGE layers and a
single frozen CodeBERT encoder, concatenated through a 128-dimensional
bottleneck---because minimality makes every design choice ablatable.
The primary novelty is not the fusion mechanism per se, but the
combination of (1) victim-contract seeding that enables authentic
multi-modal labelling, (2) enriched eight-dimensional node features
covering both degree and value-flow statistics, and (3) a full
seven-model ablation that independently justifies each choice.
